%======================================================================
\section{Conclusion}
\begin{enumerate}
\item \textbf{Timing is a risk tool.} Against the S\&P 500 proxy it roughly halves the worst drawdown ($-$50\% to between $-$21\% and $-$29\%; $-$32\% to $-$42\% at equal volatility), lost less in five of six crises, raises the Sortino and Calmar ratios and the certainty equivalent, and makes the portfolio's own risk steadier; only the steadier risk is statistically significant.
\item \textbf{It is not a return engine.} The Sharpe gain is $+$0.08 to $+$0.10, comes from 1998--2015 and is not significant in any period; terminal wealth is lower than buy and hold's because the timed portfolios hold less risk.
\item \textbf{A better forecast does not give a better portfolio.} The accuracy ranking of the five forecasts is close to the reverse of their Sharpe ranking: the least accurate, last month's RV, gives the highest Sharpe ratio and the most accurate, \VIX, the lowest. Mean timing falls from 1.01 to 0.90 as forecasts become more informative, while variance timing stays between 1.21 and 1.29. The GARCH-LSTM is barely more accurate than its GJR base and gives the same Sharpe ratio.

\textit{Why.} Variance is predictable one month ahead, returns are not, so the rule assumes a constant expected return and holds little of the market in the high-return recovery months.
\item \textbf{Training for the decision does not fix it.} With the rule, its cap, its costs and the January reset of $c$ inside the loss, the decision-trained GARCH-LSTM does worse than the accuracy-trained network on its own objective on average and on the Sharpe ratio on 4 of 6 factors, with forecasts 1.3 to 5.8 times worse, and no difference survives the multiple-testing adjustment.
\end{enumerate}
\paragraph{The open question is return information.} The best market rule we found, Xu's switches (Appendix~\ref{app:xu}), works through a 12-month trend filter, a signal about expected returns rather than a better volatility forecast. A decision-focused model would need inputs that carry information about $\mu_t$, not only about $\sigma_t^2$, to improve mean timing.

\paragraph{Limitations.} The data are revised history; VIX before 2003 is back-calculated; drawdowns are month-end; leverage is financed without a spread; equal-volatility figures are diagnostics. The 1998 start, the cap and the rules compared were settled during earlier exploration of the same data, so the multiple-testing adjustments cover the reported comparisons, not that exploration. GJR-GARCH and the GARCH-LSTM keep parameters fitted before a five-year validation block, while GARCH(1,1) is refitted every year. Decision training is one design: one-day forecasts carried 21 days forward, one value of $\gamma$, and an in-sample versus out-of-sample comparison of the objective that we have not run.
